\section{Proof sketch for Proposition~\ref{prop:ident}}
\label{app:proof}
Let atom $a$ have Cartesian position $p \in \mathbb{R}^3$ and let $v, w$ be two views with linearly
independent directions $d_v, d_w$. Its images are $q = P_v p$ and $q' = P_w p$. The back-projections are
the affine lines $\ell_v = \{x : P_v x = q\}$ (direction $d_v$) and $\ell_w = \{x : P_w x = q'\}$
(direction $d_w$). Both contain $p$; since $d_v \not\parallel d_w$, two non-parallel lines in
$\mathbb{R}^3$ that intersect do so in exactly one point, so $\ell_v \cap \ell_w = \{p\}$ and the
intersection step of Algorithm~\ref{alg:oracle} recovers $p$ exactly from the correct pairing. It remains
to show no incorrect pairing survives. A spurious pairing $(q, q'')$ with $q''$ the image of a different
atom $b$ either yields no intersection (skew lines, rejected) or yields a point $\hat{p} \neq p$ whose
projections must, by hypothesis (ii), fall outside the matching tolerance of any same-species point in at
least one remaining view, so the verification step rejects it. The verified reconstruction is
therefore exactly the atom set, on which $\spg$ returns $y(s)$ by the definition of the label.
Condition (ii) is where the argument is not vacuous: it fails precisely on projective coincidence, when
distinct atoms project within tolerance of one another in the verifying views. The measured ceiling gap
(0.0476 on the original sample) and its growth with atoms per cell (Section~\ref{sec:models}) are
measurements of that failure set.
